% WHAT A HELD-BACK NAME IS WORTH IN A TEST.
%
% \if and \ifcat compare what their two operands MEAN. A name \noexpand
% held back means \relax -- and a test reads it as NO CHARACTER AT ALL,
% which is a character code of 256 and a category of 16, so two of them
% match each other and nothing else.
%
% EXCEPT where the held-back name is an ACTIVE CHARACTER: that stands for
% the character itself, code and category 13 alike. Measured against pdftex
% with \catcode`\~=13 \def~{y}:
%
%   \if~\noexpand~                     F   (the first ~ expands to y)
%   \ifcat\noexpand~\noexpand~         T
%   \if\noexpand~\noexpand~            T
%   \ifcat\noexpand~ A                 F
%   \if\string~\noexpand~              T   <- a `~' of category 12 matches
%   \ifcat\noexpand\one\noexpand\two   T   (both are no character)
%   \if\noexpand\one\noexpand~         F   <- and one of them is not
%
% This engine read every held-back name as no character, so the last two
% came out the other way about. trip line 436 writes
% `\if$\expandafter\noexpand\dol\fi' with \dol a macro for an active $.
\catcode`\{=1 \catcode`\}=2
\tracingonline=1
\catcode`\~=13 \def~{y}
\message{[A \if~\noexpand~ T\else F\fi]}
\message{[B \ifcat\noexpand~\noexpand~ T\else F\fi]}
\message{[C \if\noexpand~\noexpand~ T\else F\fi]}
\message{[D \ifcat\noexpand~ A T\else F\fi]}
\message{[E \if\string~\noexpand~ T\else F\fi]}
\message{[F \ifcat\noexpand\undefinedone\noexpand\undefinedtwo T\else F\fi]}
\message{[G \if\noexpand\undefinedone\noexpand~ T\else F\fi]}
\message{[done]}
\end
